\chapter{Asamblare}


\section{Matrice LED}

Mini PC-ul vine echipat cu o matrice LED \ref{fig:LED-Matrix} compusă din 220 LED-uri adresabile RGB cu driver WS2812C, acestea au nevoie doar de un pin de input și sunt legate în serie, acestea formând o matrice 44x5.
Acestea sunt controlate de un Microcontroller ESP32-C3 care comunică cu modulul LattePanda prin USB, iar microcontrollerul fiind integrat direct pe \ac{PCB}-ul principal.

\begin{figure}
    \centering
    \includegraphics[width=\linewidth]{led-matrix.png}
    \caption{Matrice Led.}
    \label{fig:LED-Matrix}
\end{figure}


\section{Carcasa}
Pentru a monta cele două \ac{PCB}-uri am proiectat o carcasă folosind AutoCAD, aceasta este formată din 2 piese principale printate 3d folosind filament PLA, un schelet ce include partea din față, lateralele și partea superioară \ref{fig:Front-case}, apoi un capac pentru partea inferioară și spate \ref{fig:back-case}. Cele 2 \ac{PCB}-uri sunt montate pe schelet, acesta incluzand și un panou de difuzie printat cu filament translucid PETG.

\begin{figure}
  \centering
  \begin{subfigure}[b]{0.45\linewidth}
    \includegraphics[width=\linewidth]{case-front.png}
     \caption{Partea frontală a carcasei}
     \label{fig:Front-case}
  \end{subfigure}
  \begin{subfigure}[b]{0.45\linewidth}
    \includegraphics[width=\linewidth]{case-back.png}
    \caption{Partea din spate a carcasei}
    \label{fig:back-case}
  \end{subfigure}
  
  \caption{Cele 2 părți ale carcasei}
  \label{fig:Case}

\end{figure}

\section{Asamblare}
Pentru a asambla componentele pe \ac{PCB} am folosit șabloane din oțel inoxidabil, tăiate la laser, aceasta fiind produsă odată cu \ac{PCB}-ul, cu aceste șabloane am aplicat o pastă de fludor și am plasat fiecare componentă pe placă, începând cu matricea LED \ref{fig:LED-Matrix-Assembly} care are 220 LED-uri, apoi continuând cu \ac{PCB}-ul principal \ref{fig:Main-PCB-Assembly}, acesta având în jur de 500 de componente, dintre care TODO

\begin{figure}
    \centering
    \includegraphics[width=\linewidth]{led-matrix-asembly.png}
    \caption{Matrice Led în timpul asamblării.}
    \label{fig:LED-Matrix-Assembly}
\end{figure}

\begin{figure}
    \centering
    \includegraphics[width=0.8\linewidth]{pcb-assembly.jpg}
    \caption{PCB-ul principal în timpul asambării în timpul asamblării.}
    \label{fig:Main-PCB-Assembly}
\end{figure}


\section{Testare și depanare}

Toate defectele prezentate în această secțiune au fost identificate în urma testării prototipului și ulterior corectate în proiectul \ac{PCB}. Din considerente de cost, varianta revizuită nu a mai fost trimisă în producție.

\subsection{Defect al circuitului de alimentare}
După asamblarea plăcii, procesul de testare a evidențiat inițial o problemă majoră: regulatorul de tensiune care convertește 20V la 5V și 3,3V nu pornea. Deoarece acesta asigură și alimentarea regulatorului de 17,5 V utilizat de modulul principal, nici acesta nu devenea operațional.

Cauza problemei a fost identificată ca fiind absența unor condensatoare necesare pentru funcționarea corectă a regulatorului \ac{LDO}, utilizat în faza inițială de pornire a circuitului intern al cipului, până la activarea convertorului DC-DC. Remedierea a constat în adăugarea unui condensator, realizată printr-o conexiune suplimentară între pinul circuitului integrat și un condensator montat ulterior pe placă (Figura \ref{fig:TPS-fix}).

\begin{figure}
  \centering
  \begin{subfigure}[b]{0.45\linewidth}
    \includegraphics[width=\linewidth]{TPS-FIX.jpg}
    \caption{Firul ce leaga unul dintre LDO-uri de un condensator.}
    \label{fig:TPS-fix}
  \end{subfigure}
  \begin{subfigure}[b]{0.45\linewidth}
    \includegraphics[width=\linewidth]{USB-FIX.jpg}
    \caption{Firul USB RX-TX TX-RX. TODO BETTER PICTURE}
    \label{fig:USB-fix}
  \end{subfigure}
  
  \caption{Cele 2 părți ale carcasei}
  \label{fig:GG}

\end{figure}

\subsection{Defecte ale protocoalelor}

Sistemul a pornit corespunzător și a permis încărcarea sistemului de operare. Testele ulterioare au evidențiat însă și alte deficiențe:
\begin{itemize}
    \item Un slot \ac{NVMe} nu funcționa, cauza fiind o lipire defectuoasă a unui pin;
    \item Controlerul Ethernet \textit{Realtek RTL8125} nu pornea, fiind necesară adăugarea unui rezistor de tip pull-up pe unul dintre pini;
    \item Timpul de inițializare în BIOS era neobișnuit de mare.
\end{itemize}

Investigațiile au arătat că problema legată de BIOS era cauzată de conectarea incorectă a liniilor USB 3.2 Gen 2 către hub-ul USB. Perechile diferențiale au fost conectate eronat (RX-RX și TX-TX), în loc de configurația corectă RX-TX și TX-RX. Această eroare a apărut din utilizarea unei topologii specifice conexiunilor între cip și port, în locul celei necesare pentru interconectarea a două cipuri. Corecția a fost realizată prin refacerea conexiunilor folosind fire (Figura \ref{fig:USB-fix}).
O altă problemă identificată a fost funcționarea incorectă a LED-ului RGB frontal, utilizat pentru indicarea stării sistemului (roșu - oprit, galben - standby, alb - pornit). LED-ul rămânea aprins în culoarea roșie imediat după alimentare, iar celelalte culori nu funcționau.
Cauza acestei deficiențe a fost utilizarea incorectă a tranzistoarelor MOSFET de tip N într-o configurație de tip high-side switching, în condițiile în care tensiunea de comandă pe poartă era comparabilă sau mai mică decât tensiunea drenă-sursă. În această situație, tranzistorul nu putea comuta corect.
Soluția corectă constă în utilizarea unei configurații de tip low-side switching, în care tensiunea de comandă (3,3 V) este aplicată între poartă și masă, asigurând funcționarea corespunzătoare a tranzistorului și controlul corect al LED-ului.


\begin{figure}
  \centering
  \begin{subfigure}[b]{0.4\linewidth}
    \includegraphics[width=\linewidth]{Low vs High Side switching.png}
    \caption{Low Side vs High Side Switch}
    \label{fig:Low-high}
  \end{subfigure}
  \begin{subfigure}[b]{0.55\linewidth}
    \includegraphics[width=\linewidth]{Low-side-fix.png}
    \caption{Diagrama reparată folosind low side switch}
    \label{fig:low-side-fix}
  \end{subfigure}
  
  \caption{fix MOSFETS}
  \label{fig:2}

\end{figure}

\subsection{Alte defecte nerezolvate}

Două defecte care nu a fost reparate pe placă dar care a fost reparate în proiect sunt:
\begin{itemize}
    \item Ventilatorul pornește la putere maximă odată cu alimentarea sistemului, pentru a remedia această problemă a fost adăugat un MOSFET care permite alimentarea sistemului doar când acesta este în starea de pornit, folosind același semnal ca LED-ul alb.
    \item Al doilea controller USB nu a fi putut fi reparat pentru a funcționa la viteze usb 3.2, acesta având punctele de lipire mai inaccesibile ca primul
\end{itemize}

\subsection{Testare în sisteme de operare}
Pentru a testa performanța sistemului am folosit atât sistemul de operare Windows 11 cât și Ubuntu 26.04. 
Am testat timpul de pornire al sistemului în ambele sisteme de operare, viteza de internetului atât pe Wi-Fi cât și pe fir, viteza maximă de transmisiune pe USB dar și viteza maximă a SSD-urile conectate (fiind limitate de conexiunea PCIE x2 respecitv x1).


\begin{table}[h]
\centering
\begin{tabular}{@{}lll@{}}
\toprule
Placeholder!! & Ubuntu  & Windows \\ \midrule
Timp de pornire & 22s     & 15s     \\
Viteza LAN      & 2,5Gbps & 2,5Gbps \\
Viteza Wi-Fi    & 837mbps & 837mbps \\
Viteza SSD PCIE & 2,5Gbps & 2,5Gbps \\
Viteza USB      & 10Gbps  & 2,5Gbps \\ \bottomrule
\end{tabular}
\end{table}



